\documentclass[aspectratio=169,10pt]{beamer}

% Shared Beamer setup for Fall 2026 course slides.
% Clean Cal Poly-inspired theme: no custom class files or uncommon packages.
\mode<presentation>{
  \usetheme{default}
  \usefonttheme{professionalfonts}
}

\definecolor{CalPolyGreen}{HTML}{154734}
\definecolor{CalPolyGold}{HTML}{BD8B13}
\definecolor{CalPolySage}{HTML}{7FA28F}
\definecolor{CalPolyMint}{HTML}{DDF1E7}
\definecolor{CalPolySand}{HTML}{A29F86}
\definecolor{CalPolyBeige}{HTML}{EFEEDE}
\definecolor{DeckInk}{HTML}{1F2933}
\definecolor{DeckMuted}{HTML}{5F6B7A}
\definecolor{DeckSoft}{HTML}{F7F7F5}

\setbeamersize{text margin left=0.55in,text margin right=0.55in}
\setbeamercolor{normal text}{fg=DeckInk,bg=white}
\setbeamercolor{structure}{fg=CalPolyGreen}
\setbeamercolor{title}{fg=white}
\setbeamercolor{frametitle}{fg=CalPolyGreen,bg=white}
\setbeamercolor{block title}{bg=CalPolySage,fg=white}
\setbeamercolor{block body}{bg=CalPolyMint,fg=black}
\setbeamercolor{alerted text}{fg=CalPolyGold}

\setbeamerfont{title}{size=\Huge,series=\bfseries}
\setbeamerfont{frametitle}{size=\Large,series=\bfseries}
\setbeamerfont{block title}{size=\normalsize,series=\bfseries}

\setbeamertemplate{navigation symbols}{}
\setbeamertemplate{itemize item}{\color{CalPolyGold}$\blacktriangleright$}
\setbeamertemplate{itemize subitem}{\color{CalPolyGreen}$\bullet$}
\setbeamertemplate{footline}{
  \hfill{\color{DeckMuted}\scriptsize\insertframenumber/\inserttotalframenumber}
  \hspace{0.18in}\vspace{0.12in}
}
\setbeamertemplate{frametitle}{
  \vspace{0.18in}
  \hspace{0.02in}\insertframetitle\par
  \vspace{0.08in}
  \noindent\makebox[\linewidth][l]{%
    {\color{CalPolyGold}\rule{0.55in}{2pt}}%
    {\color{CalPolyGreen}\rule{\dimexpr\linewidth-0.55in\relax}{2pt}}%
  }
}

\newcommand{\CourseNumber}{Course}
\newcommand{\CourseTitle}{Course Title}
\newcommand{\LectureNumber}{0}
\newcommand{\LectureTitle}{Lecture Title}
\newcommand{\LectureDate}{Date}
\newcommand{\CourseUnit}{Unit}

\newcommand{\coursenumber}[1]{\renewcommand{\CourseNumber}{#1}}
\newcommand{\coursetitle}[1]{\renewcommand{\CourseTitle}{#1}}
\newcommand{\lecturenumber}[1]{\renewcommand{\LectureNumber}{#1}}
\newcommand{\lecturetitle}[1]{\renewcommand{\LectureTitle}{#1}}
\newcommand{\lecturedate}[1]{\renewcommand{\LectureDate}{#1}}
\newcommand{\courseunit}[1]{\renewcommand{\CourseUnit}{#1}}

\author{Kirk Duran}
\institute{California Polytechnic State University, San Luis Obispo}
\date{\LectureDate}

\newcommand{\makedecktitle}{
  \begingroup
  \setbeamercolor{background canvas}{bg=CalPolyGreen}
  \begin{frame}[plain]
    \vfill
    \begin{center}
      {\usebeamerfont{title}\color{white}\LectureTitle\par}
      \vspace{0.18in}
      {\Large\color{white}Lecture \LectureNumber\par}
    \end{center}
    \vfill
    \noindent{\color{white}\rule{\linewidth}{1.2pt}}\par
    \vspace{0.08in}
    \noindent\colorbox{CalPolyGold}{\makebox[0.24\linewidth][c]{\color{white}\Large\CourseNumber}}
    \hspace{0.12in}{\color{white}\Large Kirk Duran}
  \end{frame}
  \endgroup
}

\newenvironment{keyidea}{\begin{block}{Key idea}}{\end{block}}
\newenvironment{activity}{\begin{block}{In class}}{\end{block}}
\newenvironment{cpdefinition}{\begin{block}{Definition}}{\end{block}}
\newenvironment{exampleblockcp}{
  \begingroup
  \setbeamercolor{block title}{bg=CalPolySand,fg=white}
  \setbeamercolor{block body}{bg=CalPolyBeige,fg=black}
  \begin{block}{Example}
}{
  \end{block}
  \endgroup
}

\newcommand{\imageplaceholder}[2][0.3in]{%
  \noindent\makebox[\linewidth][c]{%
    \fcolorbox{CalPolySage}{CalPolyBeige}{%
      \parbox[c][#1][c]{0.86\linewidth}{%
        \centering
        {\color{DeckMuted}\scriptsize Image placeholder: }%
        {\color{DeckInk}\scriptsize\ttfamily\detokenize{#2}}%
      }%
    }%
  }%
}

\newcommand{\sectionframe}[1]{
  \begingroup
  \setbeamercolor{background canvas}{bg=CalPolyGreen}
  \begin{frame}[plain]
    \vfill
    {\color{white}\LARGE\bfseries #1\par}
    \vspace{0.18in}
    {\color{CalPolyGold}\rule{0.22\paperwidth}{3pt}}
    {\color{white}\rule{0.62\paperwidth}{3pt}}
    \vfill
  \end{frame}
  \endgroup
}

\newcommand{\placeholdernote}{
  \begin{block}{Placeholder}
    Replace these bullets with the final lecture material, examples, and in-class activities.
  \end{block}
}


\pdfstringdefDisableCommands{\def\translate#1{#1}}

\graphicspath{{../assets/lecture-24/}}

% If an image file is not yet available, the slide displays a labeled
% placeholder using the intended filename. Place the matching file in
% ../assets/lecture-24/ to replace the placeholder automatically.
\newcommand{\lectureimage}[2][1.75in]{%
  \IfFileExists{../assets/lecture-24/#2}{%
    \includegraphics[width=\linewidth,height=#1,keepaspectratio]{#2}%
  }{%
    \fbox{%
      \parbox[c][#1][c]{0.94\linewidth}{%
        \centering
        \small\textbf{Image Placeholder}\\[0.08in]
        \scriptsize\texttt{\detokenize{#2}}
      }%
    }%
  }%
}

\setbeamersize{text margin left=0.42in,text margin right=0.42in}

\coursenumber{CS 1001}
\coursetitle{Introduction to Computing}
\lecturenumber{24}
\lecturetitle{Higher-Level Sequence Patterns}
\lecturedate{Fall 2026}
\courseunit{9. Iteration}

\begin{document}

\makedecktitle

\begin{frame}{Today}
  \begin{itemize}
    \item Move from reasoning about \textbf{loop mechanics} to reasoning about \textbf{sequence-level intent}.
    \item Give formal names to transformation and selection patterns already implemented with loops.
    \item Refactor suitable loop-based list construction into list comprehensions.
    \item Trace comprehensions using the same binding, condition, and evaluation model used for \texttt{for} loops.
    \item Distinguish computations that fit a comprehension from those that remain clearer as explicit loops.
    \item Look briefly at the runtime mechanism that allows \texttt{for} to traverse a sequence.
  \end{itemize}

  \begin{keyidea}
    Lecture 24 does not introduce a new kind of computation. It introduces a \textbf{higher-level representation} for some computations students already know how to express with loops.
  \end{keyidea}
\end{frame}

\begin{frame}{The Boundary for This Lecture}
  \begin{columns}[T,onlytextwidth]
    \begin{column}{0.47\textwidth}
      \begin{block}{Today}
        \begin{itemize}
          \item map / transformation as a pattern
          \item filter / selection as a pattern
          \item list comprehensions
          \item semantic equivalence with simple loops
          \item choosing a readable representation
        \end{itemize}
      \end{block}
    \end{column}
    \begin{column}{0.47\textwidth}
      \begin{alertblock}{Not today}
        \begin{itemize}
          \item \texttt{map()} or \texttt{filter()} built-ins
          \item a catalog of sequence built-ins
          \item nested comprehensions
          \item nested data structures
          \item searching or sorting algorithms
        \end{itemize}
      \end{alertblock}
    \end{column}
  \end{columns}
\end{frame}

\sectionframe{Part 1: From Loop Mechanics to Computational Intent}

\begin{frame}{The Iteration Unit Has Three Levels}
  \begin{center}
    \renewcommand{\arraystretch}{1.35}
    \begin{tabular}{l|l|l}
      \textbf{Lecture} & \textbf{Primary question} & \textbf{Level} \\\hline
      22 & How does a \texttt{for} loop execute? & mechanism \\
      23 & What recurring loop structures solve common tasks? & construction \\
      24 & What computation does this loop represent? & abstraction
    \end{tabular}
  \end{center}

  \begin{keyidea}
    We are moving upward in abstraction, not replacing the execution model from Lecture 22.
  \end{keyidea}
\end{frame}

\begin{frame}[t]{From Mechanism to Intent}
  \small
  \begin{columns}[T,onlytextwidth]
    \begin{column}{0.54\textwidth}
      \begin{block}{Mechanistic description}
        Create an empty list, visit each input value, compute a result, append that result, repeat until the input is exhausted.
      \end{block}

      \begin{block}{Intent-level description}
        Produce one transformed output value for each input value.
      \end{block}

      \begin{keyidea}
        Both descriptions are correct. They answer different questions.
      \end{keyidea}
    \end{column}
    \begin{column}{0.40\textwidth}
      % Intended visual: ladder from statement-level loop mechanics to sequence-level intent.
      \lectureimage[2.25in]{abstraction-ladder.png}
    \end{column}
  \end{columns}
\end{frame}

\begin{frame}{Formal Names for Familiar Patterns}
  \begin{center}
    \renewcommand{\arraystretch}{1.35}
    \begin{tabular}{l|l|l}
      \textbf{Pattern} & \textbf{Relationship} & \textbf{Typical result} \\\hline
      transformation / map & each input produces a related output & new sequence \\
      selection / filter & only qualifying inputs remain & new sequence \\
      accumulation & repeated updates summarize the input & one value
    \end{tabular}
  \end{center}

  \begin{alertblock}{Terminology}
    \textbf{Map} and \textbf{filter} name computational patterns here. We are not introducing Python's \texttt{map()} or \texttt{filter()} functions.
  \end{alertblock}

  Naming the pattern lets us recognize the same computational relationship even when the particular arithmetic or data changes.
\end{frame}

\begin{frame}{Refactoring Should Preserve Meaning}
  \begin{cpdefinition}
    A \textbf{refactoring} changes how code is expressed while preserving its intended observable behavior.
  \end{cpdefinition}

  When we replace a simple list-building loop with a comprehension, we should preserve:

  \begin{itemize}
    \item which input values are examined;
    \item which conditions are checked;
    \item which result values are produced;
    \item the order of the produced values; and
    \item whether the original input is mutated.
  \end{itemize}

  \begin{keyidea}
    Shorter syntax is useful only when it preserves both behavior and clarity.
  \end{keyidea}
\end{frame}

\sectionframe{Part 2: List Comprehensions as Refactored List Construction}

\begin{frame}{The Same Relationship as a List Comprehension}
  \begin{columns}[T,onlytextwidth]
    \begin{column}{0.51\textwidth}
      \begin{block}{Explicit construction}
        \texttt{result = []}\\
        \texttt{for value in values:}\\
        \hspace{0.18in}\texttt{result.append(value * 2)}
      \end{block}
    \end{column}
    \begin{column}{0.43\textwidth}
      \begin{block}{Comprehension}
        \texttt{result = [value * 2}\\
        \hspace{0.45in}\texttt{for value in values]}
      \end{block}
    \end{column}
  \end{columns}

  \begin{keyidea}
    The comprehension removes bookkeeping that is predictable from the pattern: creating the empty list and explicitly calling \texttt{append}.
  \end{keyidea}
\end{frame}

\begin{frame}{A List Comprehension Is an Expression}
  \begin{cpdefinition}
    A \textbf{list comprehension} is an expression that constructs a new list by repeatedly evaluating an output expression over values obtained from an input sequence.
  \end{cpdefinition}

  \begin{block}{General form}
    \centering
    \texttt{[output\_expression for name in sequence]}
  \end{block}

  Because it is an expression, the completed list can be:
  \begin{itemize}
    \item assigned to a name;
    \item returned from a function; or
    \item used as part of another expression.
  \end{itemize}
\end{frame}

\begin{frame}{Anatomy of a Comprehension}
  \begin{columns}[T,onlytextwidth]
    \begin{column}{0.55\textwidth}
      \begin{block}{Example}
        \centering
        \texttt{[x + 1 for x in values]}
      \end{block}

      \begin{center}
        \renewcommand{\arraystretch}{1.28}
        \begin{tabular}{l|l}
          \textbf{Part} & \textbf{Role} \\\hline
          \texttt{values} & source sequence \\
          \texttt{x} & repeatedly bound name \\
          \texttt{x + 1} & output expression
        \end{tabular}
      \end{center}
    \end{column}
    \begin{column}{0.39\textwidth}
      % Intended visual: input values flowing through a transform box into a new list.
      \lectureimage[2.45in]{comprehension-pipeline.png}
    \end{column}
  \end{columns}
\end{frame}

\begin{frame}[t]{Trace It Using the Old Loop Model}
  \small
  For:
  \begin{block}{Expression}
    \centering
    \texttt{[x + 1 for x in [1, 2, 3, 4]]}
  \end{block}

  \begin{center}
    \renewcommand{\arraystretch}{1.15}
    \begin{tabular}{c|c|c}
      \textbf{Step} & \textbf{x} & \textbf{Produced value} \\\hline
      1 & 1 & 2 \\
      2 & 2 & 3 \\
      3 & 3 & 4 \\
      4 & 4 & 5
    \end{tabular}
  \end{center}

  Final value: \texttt{[2, 3, 4, 5]}

  \begin{keyidea}
    The compact syntax does not change the binding model: obtain a value, bind the name, evaluate the output expression, and add the produced value to the new list.
  \end{keyidea}
\end{frame}

\begin{frame}{The Output Expression Can Call a Function}
  \begin{block}{Code}
    \texttt{def square(n: int) -> int:}\\
    \hspace{0.20in}\texttt{return n * n}\\[0.10in]
    \texttt{squares = [square(x) for x in [1, 2, 3, 4]]}
  \end{block}

  The comprehension causes four separate calls:

  \begin{center}
    \texttt{square(1)} $\rightarrow$ \texttt{1}\qquad
    \texttt{square(2)} $\rightarrow$ \texttt{4}\qquad
    \texttt{square(3)} $\rightarrow$ \texttt{9}\qquad
    \texttt{square(4)} $\rightarrow$ \texttt{16}
  \end{center}
\end{frame}

\begin{frame}{The Input Is Not the Output}
  \begin{columns}[T,onlytextwidth]
    \begin{column}{0.46\textwidth}
      \begin{block}{Input}
        \texttt{values = [2, 4, 6]}
      \end{block}
    \end{column}
    \begin{column}{0.46\textwidth}
      \begin{block}{Comprehension}
        \texttt{doubled = [x * 2 for x in values]}
      \end{block}
    \end{column}
  \end{columns}

  Afterwards:
  \begin{center}
    \texttt{values} $\rightarrow$ \texttt{[2, 4, 6]}\\
    \texttt{doubled} $\rightarrow$ \texttt{[4, 8, 12]}
  \end{center}

  \begin{keyidea}
    A list comprehension constructs a \textbf{new list}. It does not implicitly mutate the source list.
  \end{keyidea}
\end{frame}

\sectionframe{Part 3: Selection and Conditional Production}

\begin{frame}{Selection Has the Same Refactoring Idea}
  \begin{columns}[T,onlytextwidth]
    \begin{column}{0.51\textwidth}
      \begin{block}{Explicit loop}
        \texttt{result = []}\\
        \texttt{for x in values:}\\
        \hspace{0.18in}\texttt{if x > 0:}\\
        \hspace{0.36in}\texttt{result.append(x)}
      \end{block}
    \end{column}
    \begin{column}{0.43\textwidth}
      \begin{block}{Filtered comprehension}
        \texttt{result = [x for x in values}\\
        \hspace{0.48in}\texttt{if x > 0]}
      \end{block}
    \end{column}
  \end{columns}

  The condition determines whether a candidate contributes a value to the result.
\end{frame}

\begin{frame}{General Form with a Condition}
  \begin{block}{Filtered comprehension}
    \centering
    \texttt{[output\_expression for name in sequence if condition]}
  \end{block}

  For each candidate:
  \begin{enumerate}
    \item bind the candidate to \texttt{name};
    \item evaluate \texttt{condition};
    \item if the condition is \texttt{True}, evaluate the output expression;
    \item otherwise produce nothing for that candidate.
  \end{enumerate}
\end{frame}

\begin{frame}{Trace Selection Explicitly}
  \begin{block}{Expression}
    \centering
    \texttt{[x for x in [-2, 5, 0, 7] if x > 0]}
  \end{block}

  \begin{center}
    \renewcommand{\arraystretch}{1.22}
    \begin{tabular}{c|c|c}
      \textbf{x} & \textbf{x > 0} & \textbf{Produced value} \\\hline
      -2 & False & --- \\
      5 & True & 5 \\
      0 & False & --- \\
      7 & True & 7
    \end{tabular}
  \end{center}

  Final value: \texttt{[5, 7]}

  Selected values appear in the same relative order in which they were encountered.
\end{frame}

\begin{frame}{Selection and Transformation Can Be Combined}
  \begin{block}{Example}
    \centering
    \texttt{[x * 2 for x in values if x > 0]}
  \end{block}

  Conceptually:
  \begin{center}
    \textbf{candidate} $\longrightarrow$ \textbf{condition} $\longrightarrow$ \textbf{transform if selected}
  \end{center}

  \begin{itemize}
    \item some candidates produce no output value;
    \item selected candidates can be transformed before entering the result;
    \item the result is still a newly constructed list.
  \end{itemize}
\end{frame}

\begin{frame}{Written Order Is Not Evaluation Order}
  Consider:
  \begin{block}{Expression}
    \centering
    \texttt{[transform(x) for x in values if check(x)]}
  \end{block}

  For each \texttt{x}, Python conceptually does:
  \begin{enumerate}
    \item evaluate \texttt{check(x)};
    \item if it is \texttt{True}, evaluate \texttt{transform(x)};
    \item place the transformed value into the result.
  \end{enumerate}

  \begin{alertblock}{Common reading mistake}
    The output expression appears first syntactically, but it is not evaluated for candidates rejected by the condition. If \texttt{check(x)} is false, \texttt{transform(x)} is not called for that candidate.
  \end{alertblock}
\end{frame}

\begin{frame}{A Reliable Mental Expansion}
  When a comprehension feels dense, expand it mentally.

  \begin{block}{Comprehension}
    \texttt{result = [f(x) for x in values if p(x)]}
  \end{block}

  \begin{block}{Equivalent reasoning structure}
    \texttt{result = []}\\
    \texttt{for x in values:}\\
    \hspace{0.18in}\texttt{if p(x):}\\
    \hspace{0.36in}\texttt{result.append(f(x))}
  \end{block}

  \begin{keyidea}
    The explicit loop remains the tracing model even when the program uses compact syntax.
  \end{keyidea}
\end{frame}

\sectionframe{Part 4: How Do Loops Work ``Under the Engine''?}

\begin{frame}[t]{The Surface Syntax Hides a Runtime Conversation}
  \small
  \begin{columns}[T,onlytextwidth]
    \begin{column}{0.54\textwidth}
      When Python executes:
      \begin{block}{Surface syntax}
        \texttt{for value in sequence:}\\
        \hspace{0.20in}\texttt{use(value)}
      \end{block}

      it needs a runtime mechanism that can answer two questions repeatedly:
      \begin{enumerate}
        \item Is there another value?
        \item If so, what is it?
      \end{enumerate}

      \begin{alertblock}{Flavor, not required syntax}
        You are not expected to write the lower-level mechanism yourself in this course.
      \end{alertblock}
    \end{column}
    \begin{column}{0.40\textwidth}
      % Intended visual: a sequence feeding values through an iterator state into a loop variable.
      \lectureimage[2.25in]{iterator-engine.png}
    \end{column}
  \end{columns}
\end{frame}

\begin{frame}{Conceptual Runtime Model}
  At a high level, Python behaves approximately like this:

  \begin{block}{Conceptual model}
    \texttt{iterator = iter(sequence)}\\
    \texttt{while another value is available:}\\
    \hspace{0.20in}\texttt{value = next(iterator)}\\
    \hspace{0.20in}\texttt{execute loop body}
  \end{block}

  Internally, the iterator maintains enough state to know which value should come next. When no next value exists, the runtime reports \texttt{StopIteration}; the \texttt{for} statement handles that termination automatically.

  \begin{keyidea}
    The \texttt{for} statement is a high-level control structure built on a lower-level mechanism for requesting successive values.
  \end{keyidea}
\end{frame}

\begin{frame}{Why \texttt{range} Works Naturally with \texttt{for}}
  A \texttt{range} value represents a numerical progression using a start, stop, and step.

  \begin{block}{Example}
    \centering
    \texttt{for i in range(2, 8, 2):}
  \end{block}

  The loop receives:
  \begin{center}
    \texttt{2} $\rightarrow$ \texttt{4} $\rightarrow$ \texttt{6}
  \end{center}

  The range value can support iteration without first constructing a list containing every integer in the progression.
\end{frame}

\begin{frame}[t]{Comprehensions Use the Same Iteration Machinery}
  \small
  A comprehension does not remove iteration from the computation.

  \begin{columns}[T,onlytextwidth]
    \begin{column}{0.52\textwidth}
      \begin{block}{Source}
        \texttt{[x * 2 for x in values if x > 0]}
      \end{block}

      Conceptually, Python still:
      \begin{itemize}
        \item obtains successive source values;
        \item performs repeated bindings;
        \item checks the condition;
        \item evaluates the output expression;
        \item constructs a new list.
      \end{itemize}
    \end{column}
    \begin{column}{0.42\textwidth}
      % Intended visual: same iterator feeding either an explicit loop body or a comprehension pipeline.
      \lectureimage[2.00in]{loop-and-comprehension-runtime.png}
    \end{column}
  \end{columns}

  The abstraction hides predictable bookkeeping while preserving the same underlying traversal.
\end{frame}

\sectionframe{Part 5: Choosing the Right Representation}

\begin{frame}{A Comprehension Has a Narrow Job}
  A comprehension is usually a strong fit when all of the following are true:
  \begin{itemize}
    \item the goal is to construct one new list;
    \item values come from one straightforward source sequence;
    \item each selected candidate produces at most one output value;
    \item the transformation is short enough to understand locally;
    \item the condition, if present, is similarly simple.
  \end{itemize}

  \begin{block}{Good fit}
    \centering
    \texttt{[word.lower() for word in words if word != ""]}
  \end{block}
\end{frame}

\begin{frame}{An Explicit Loop Is Better for Richer State Changes}
  Prefer a normal loop when the body coordinates several pieces of state or several actions.

  \begin{itemize}
    \item multiple accumulators;
    \item several independent updates;
    \item complex conditional branching;
    \item intentional mutation of existing data;
    \item multiple side effects;
    \item intermediate states that matter for debugging or explanation.
  \end{itemize}

  \begin{keyidea}
    A comprehension is not an ``advanced loop.'' It is a specialized expression for constructing a list.
  \end{keyidea}
\end{frame}

\begin{frame}[t]{Shorter Is Not Automatically Clearer}
  \small
  \begin{columns}[T,onlytextwidth]
    \begin{column}{0.48\textwidth}
      \begin{block}{Compact and clear}
        \texttt{[x * 2 for x in values if x > 0]}
      \end{block}
    \end{column}
    \begin{column}{0.48\textwidth}
      \begin{block}{Too much hidden in one expression}
        several conditions + several calls + several transformations
      \end{block}
    \end{column}
  \end{columns}

  \begin{center}
    \lectureimage[1.05in]{readability-spectrum.png}
  \end{center}

  \begin{keyidea}
    Choose the representation that makes the computation easiest to verify, explain, and maintain.
  \end{keyidea}
\end{frame}

\begin{frame}{Debugging Can Favor the Explicit Form}
  Suppose a comprehension produces the wrong list.

  One useful debugging move is to temporarily expand it:

  \begin{block}{Compact version}
    \texttt{result = [f(x) for x in values if p(x)]}
  \end{block}

  \begin{block}{Debugging version}
    \texttt{result = []}\\
    \texttt{for x in values:}\\
    \hspace{0.18in}\texttt{if p(x):}\\
    \hspace{0.36in}\texttt{result.append(f(x))}
  \end{block}

  The expanded form exposes intermediate state and gives the debugger more obvious stopping points.
\end{frame}

\begin{frame}{Semantic Equivalence Checklist}
  Before replacing a loop with a comprehension, verify:

  \begin{enumerate}
    \item Is the same source sequence traversed?
    \item Are candidates examined in the same order?
    \item Is the same condition applied to each candidate?
    \item Is the same transformation applied to selected candidates?
    \item Is a fresh result list produced?
    \item Is the original input left unchanged when that is the intended contract?
  \end{enumerate}
\end{frame}

\sectionframe{Part 6: Synthesis}

\begin{frame}{One Problem, Three Levels of Description}
  Consider producing the lengths of every word.

  \begin{center}
    \renewcommand{\arraystretch}{1.35}
    \begin{tabular}{l|l}
      \textbf{Level} & \textbf{Description} \\\hline
      execution & repeatedly bind a word and execute statements \\
      loop pattern & build a new list by transforming each input \\
      sequence pattern & map each word to its length
    \end{tabular}
  \end{center}

  \begin{block}{Comprehension representation}
    \centering
    \texttt{[len(word) for word in words]}
  \end{block}
\end{frame}

\begin{frame}{Map, Filter, and Accumulate Are Different Relationships}
  \begin{center}
    \renewcommand{\arraystretch}{1.3}
    \begin{tabular}{l|l|l}
      \textbf{Pattern} & \textbf{Input-to-output relationship} & \textbf{Natural representation} \\\hline
      map & one output per input & loop or comprehension \\
      filter & zero or one output per input & loop or comprehension \\
      accumulate & many inputs update shared summary state & explicit loop
    \end{tabular}
  \end{center}

  \begin{keyidea}
    Not every repeated computation should become a comprehension.
  \end{keyidea}
\end{frame}

\begin{frame}{Integrated Prediction}
  Do not treat this as one opaque expression.

  \begin{block}{Code}
    \texttt{def adjust(n: int) -> int:}\\
    \hspace{0.20in}\texttt{return n * 3}\\[0.08in]
    \texttt{values = [-2, 0, 1, 4]}\\
    \texttt{result = [adjust(x) for x in values if x > 0]}
  \end{block}

  Predict:
  \begin{enumerate}
    \item every value bound to \texttt{x};
    \item every result of \texttt{x > 0};
    \item exactly which calls to \texttt{adjust} occur;
    \item the final value of \texttt{result}.
  \end{enumerate}
\end{frame}

\begin{frame}{Integrated Trace}
  \begin{center}
    \renewcommand{\arraystretch}{1.25}
    \begin{tabular}{c|c|c}
      \textbf{x} & \textbf{x > 0} & \textbf{adjust(x)} \\\hline
      -2 & False & not called \\
      0 & False & not called \\
      1 & True & 3 \\
      4 & True & 12
    \end{tabular}
  \end{center}

  Therefore:
  \begin{center}
    \texttt{result = [3, 12]}
  \end{center}

  The compressed syntax still obeys the same rules for bindings, conditions, function calls, and returned values.
\end{frame}

\begin{frame}{Iteration Unit: The Complete Picture}
  \begin{center}
    \renewcommand{\arraystretch}{1.32}
    \begin{tabular}{l|l}
      \textbf{Lecture} & \textbf{Core responsibility} \\\hline
      22 & execution mechanics of \texttt{for} and \texttt{range} \\
      23 & constructing useful computations with recurring loop patterns \\
      24 & recognizing intent and selecting a higher-level representation
    \end{tabular}
  \end{center}

  \begin{keyidea}
    The strongest mental model works at multiple levels: understand the mechanism, recognize the pattern, then choose the clearest representation.
  \end{keyidea}
\end{frame}

\begin{frame}{Takeaways}
  \begin{enumerate}
    \item Map and filter name relationships between an input sequence and a result sequence.
    \item A list comprehension is an expression for constructing a new list from those relationships.
    \item Comprehensions still rely on repeated binding and iteration; they hide bookkeeping rather than eliminate computation.
    \item Conditions are evaluated before the output expression for a candidate is produced.
    \item Explicit loops remain preferable when the computation manages richer state or several actions.
    \item Understanding the lower-level iteration mechanism explains why the same \texttt{for} syntax works across different iterable values.
  \end{enumerate}

  \begin{block}{Next}
    We leave sequential data as the primary structure and introduce \textbf{dictionaries and associative data}.
  \end{block}
\end{frame}

\end{document}
